\documentclass[a4paper,10pt]{report}\input{../0.Préambule/Préambule_cours_prof}\begin{document}

\newpage
\section*{Statistiques}
\addcontentsline{toc}{section}{Statistiques}

\vspace{-3mm}
\subsection*{Effectifs et fréquences}
\addcontentsline{toc}{subsection}{Effectifs et fréquences}

\vspace{-3mm}
\begin{exer1}
On a écrit la même expression dans différentes langues (néerlandaise, italienne, anglaise, allemande et française). Calcule l'effectif, puis la fréquence des voyelles dans chaque expression.

\noindent \begin{center}
Gelukkige verjaardag \hspace{3mm} Buon compleanno \hspace{3mm} Happy birthday \hspace{3mm} Alles Gute zum Geburtstag \hspace{3mm} Joyeux anniversaire
\end{center}
\end{exer1}

\begin{exer1}
Dans les classes de 5$^e$ TANDEM et de 5$^e$ COUBERTIN, $32$ élèves sont demi-pensionnaires dont $14$ sont en 5$^e$ COUBERTIN. Les $11$ autres élèves de la classe sont externes \\comme $9$ élèves de la classe de 5$^e$ TANDEM.

\begin{enumerate}[font=\color{exer1}]
\item Complète le tableau suivant
\item Combien d'élèves y a-t-il en 5$^e$ TANDEM.
\item Quel est le nombre total d'élèves pour ces deux classes
\end{enumerate}

\vspace{-3cm}
\begin{flushright}
\renewcommand{\arraystretch}{1.5}
\begin{tabularx}{6cm}{|p{1.5cm}|*{3}{>{\centering \arraybackslash}X|}}
\cline{2-4}
\multicolumn{1}{c|}{} & \cellcolor{exer1!20}{5$^e$ T} & \cellcolor{exer1!20}{5$^e$ C} & \cellcolor{exer1!20}{Total} \\\hline
\cellcolor{exer1!20}{Externes} & & &\\\hline
\cellcolor{exer1!20}{D.P.} & & &\\\hline
\cellcolor{exer1!20}{Total} & & &\\\hline
\end{tabularx}
\end{flushright}
\end{exer1}

\begin{exer1}
Le diagramme en bâton ci-dessous représente le temps de trajet pour venir au collège des élèves de 5$^e$ TANDEM.

\noindent \begin{minipage}{8cm}
\begin{enumerate}[font=\color{exer1}]
\item Calcule l'effectif total de la classe de 5$^e$ TANDEM.
\item Quelle est la fréquence des élèves dont le temps de trajet est égale à $15$ minutes ?
\item Quelle est la fréquence dont le temps de trajet est égale à $45$ minutes ?
\item Quelle est la fréquence des élèves dans le temps de trajet est compris entre $5$ et $30$ minutes incluses ? Écris le résultats sous la forme d'un pourcentage.
\end{enumerate}
\end{minipage}
\hspace{5mm}
\begin{minipage}{8cm}
\begin{center}
\begin{tikzpicture}[line cap=round,line join=round,>=triangle 45,x=1.0cm,y=1.0cm]
\begin{axis}[x=0.15cm,y=0.6cm,axis lines=middle,ymajorgrids=true,xmajorgrids=true,
xmin=-0.8,xmax=50,ymin=-0.15,ymax=6.6,
xtick={-0.0,5.0,...,45},ytick={-0,1,...,6},
xlabel=Temps (en mn),
ylabel=Effectif,]
\clip(-1.45,-0.2) rectangle (50,6.6);
\draw[line width=0.8pt,color=exer1,fill=exer1,fill opacity=0.3] (4,0.) rectangle (6,4.);
\draw[line width=0.8pt,color=exer1,fill=exer1,fill opacity=0.3] (9,0.) rectangle (11,5.);
\draw[line width=0.8pt,color=exer1,fill=exer1,fill opacity=0.3] (14,0.) rectangle (16,6.);
\draw[line width=0.8pt,color=exer1,fill=exer1,fill opacity=0.3] (19,0.) rectangle (21,3.);
\draw[line width=0.8pt,color=exer1,fill=exer1,fill opacity=0.3] (24,0.) rectangle (26,2.);
\draw[line width=0.8pt,color=exer1,fill=exer1,fill opacity=0.3] (29,0.) rectangle (31,4.);
\draw[line width=0.8pt,color=exer1,fill=exer1,fill opacity=0.3] (34,0.) rectangle (36,2.);
\draw[line width=0.8pt,color=exer1,fill=exer1,fill opacity=0.3] (39,0.) rectangle (41,2.);
\draw[line width=0.8pt,color=exer1,fill=exer1,fill opacity=0.3] (44,0.) rectangle (46,2.);
\end{axis}
\end{tikzpicture}
\end{center}
\end{minipage}
\end{exer1}

\begin{exer1}
Ces tableaux donnent la répartition des masses des œufs (en grammes) d'un élevage de poules.

\begin{center}
\renewcommand{\arraystretch}{1.5}
\begin{tabularx}{\linewidth}{|p{2cm}|*{17}{>{\centering \arraybackslash}X|}}
\hline
\cellcolor{exer1!20}{Masse (en g)}  & $41$- & $42$ & $43$ & $44$ & $45$ & $46$ & $47$ & $48$ & $49$ & $50$ & $51$ & $52$ & $53$ & $54$ & $55$ & $56$ & $57$ \\\hline
\cellcolor{exer1!20}{Effectif} & $4$ & $1$ & $2$ & $1$ & $2$ & $2$ & $3$ & $2$ & $2$ & $3$ & $4$ & $4$ & $10$ & $15$ & $17$ & $30$ & $46$ \\\hline
\end{tabularx}
\end{center}

\begin{center}
\renewcommand{\arraystretch}{1.5}
\begin{tabularx}{\linewidth}{|p{2cm}|*{16}{>{\centering \arraybackslash}X|}}
\hline
\cellcolor{exer1!20}{Masse (en g)} & $58$ & $59$ & $60$ & $61$ & $62$ & $63$ & $64$ & $65$ & $66$ & $67$ & $68$ & $69$ & $70$ & $71$ & $72$ & $73$+ \\\hline
\cellcolor{exer1!20}{Effectif} & $39$ & $48$ & $57$ & $55$ & $53$ & $68$ & $72$ & $91$ & $94$ & $93$ & $85$ & $75$ & $68$ & $59$ & $55$ & $140$ \\\hline
\end{tabularx}
\end{center}

\begin{enumerate}[font=\color{exer1}]
\item Calcule le nombre d'œufs répertorié.
\item Suivant le pays les œufs ne sont pas calibré de la même façon. \\Complète la colonne $E$ (Effectif) \\pour chaque tableau.

\vspace{-8mm}
\noindent \renewcommand{\arraystretch}{1.5}
\begin{tabularx}{5cm}{|p{1cm}|p{1.5cm}|*{2}{>{\centering \arraybackslash}X|}}
\hline
\multicolumn{4}{|c|}{\cellcolor{exer1!20}{Suisse}}\\\hline
\multicolumn{2}{|c|}{Calibre} & E & F\\\hline
\cellcolor{exer1!20}{Petit} & $49$g et - & &  \\\hline
\cellcolor{exer1!20}{Moyen} & $50$g à $65$g & &  \\\hline
\cellcolor{exer1!20}{Gros} & $66$g et + & &  \\\hline
\end{tabularx}
\renewcommand{\arraystretch}{1.5}
\begin{tabularx}{4.5cm}{|p{0.5cm}|p{1.5cm}|*{2}{>{\centering \arraybackslash}X|}}
\hline
\multicolumn{4}{|c|}{\cellcolor{exer1!20}{France}}\\\hline
\multicolumn{2}{|c|}{Calibre} & E & F\\\hline
\cellcolor{exer1!20}{S} & $52$g et - & &  \\\hline
\cellcolor{exer1!20}{M} & $53$g à $62$g & &  \\\hline
\cellcolor{exer1!20}{L} & $62$g à $72$g & &  \\\hline
\cellcolor{exer1!20}{XL} & $73$g et + & &  \\\hline
\end{tabularx}
\renewcommand{\arraystretch}{1.5}
\begin{tabularx}{5.5cm}{|p{1.5cm}|p{1.5cm}|*{2}{>{\centering \arraybackslash}X|}}
\hline
\multicolumn{4}{|c|}{\cellcolor{exer1!20}{Canada}}\\\hline
\multicolumn{2}{|c|}{Calibre} & E & F\\\hline
\cellcolor{exer1!20}{Pee Wee} & $41$g et - & &  \\\hline
\cellcolor{exer1!20}{Petit} & $42$g à $48$g & &  \\\hline
\cellcolor{exer1!20}{Moyen} & $49$g à $55$g & &  \\\hline
\cellcolor{exer1!20}{Gros} & $56$g à $63$g & &  \\\hline
\cellcolor{exer1!20}{Extra Gros} & $64$g à $69$g & &  \\\hline
\cellcolor{exer1!20}{Jumbo} & $70$g et + & &  \\\hline
\end{tabularx}

\vspace{-8mm}
\item Complète la colonne $F$ \\(fréquence en pourcentage) de chaque tableau.
\item Compare le pourcentages obtenus dans chaque pays pour la catégorie Gros (L pour la France).
\end{enumerate}
\end{exer1}

\begin{exer1}
Dans le collège, on a demandé au élèves de 5$^e$, le sport qu'ils pratiquent au sein d'un club. Le tableau suivant récapitule les différentes réponses. Complète le.

\begin{center}
\renewcommand{\arraystretch}{1.5}
\begin{tabularx}{\linewidth}{|p{2cm}|*{7}{>{\centering \arraybackslash}X|}}
\cline{2-4}
\rowcolor{exer1!20}Sport  & Football & Basket & Rugby & Tennis & Danse & Equitation & Total \\\hline
\cellcolor{exer1!20}{Effectif} & $52$ & $48$ & $35$ & $27$ & $46$ & & $250$\\\hline
\cellcolor{exer1!20}{Fréquence} & & & & & & & \\\hline
\end{tabularx}
\end{center}
\end{exer1}

\begin{exer2}
La tableau représente le choix de la langue vivante en classe de cinquième. Des données ont été effacées. Retrouve les.

\vspace{-5mm}
\begin{center}
\renewcommand{\arraystretch}{1.5}
\begin{tabularx}{\linewidth}{|p{3.5cm}|*{4}{>{\centering \arraybackslash}X|}}
\cline{2-5}
\multicolumn{1}{c|}{} & \cellcolor{exer2!20}{Anglais} & \cellcolor{exer2!20}{Espagnol} & \cellcolor{exer2!20}{Allemand} & \cellcolor{exer2!20}{Total} \\\hline
\cellcolor{exer2!20}{Garçons} & & & $5$ & $58$\\\hline
\cellcolor{exer2!20}{Filles} & $38$ & $15$ & &\\\hline
\cellcolor{exer2!20}{Total} & $82$ & & & $122$\\\hline
\end{tabularx}
\end{center}
\end{exer2}

\subsection*{Moyenne une série statistique}
\addcontentsline{toc}{subsection}{Moyenne une série statistique}

\begin{exer1}
Les élèves de $5^e$ du collège de Yenne ont indiqué le nombre de livres qu'ils ont lus durant le mois de septembre. Voici les résultats de l'enquête.

\begin{center}
\renewcommand{\arraystretch}{1.5}
\begin{tabularx}{\linewidth}{|p{3cm}|*{7}{>{\centering \arraybackslash}X|}}
\hline
\cellcolor{exer1!20}{Nombre de livre lus} & $0$ & $1$ & $2$ & $3$ & $7$ & $8$ & $15$\\\hline
\cellcolor{exer1!20}{Effectifs} & $12$ & $4$ & $3$ & $3$ & $1$ & $1$ & $1$ \\\hline
\end{tabularx}
\end{center}

Calcul le nombre de livres lus, en moyenne, pas un élève du collège de Yenne au mois de septembre.
\end{exer1}

\begin{exer1}
Une équipe de volley comporte neuf joueurs. Voici leurs tailles et le nombre de points que chacun a marqués lors de cette saison.

\begin{center}
\renewcommand{\arraystretch}{1.5}
\begin{tabularx}{\linewidth}{|*{9}{>{\centering \arraybackslash}X|}}
\hline
\rowcolor{exer1!20} Marc & Akim & Alex & Loic & Chris & Olivier & Sylvain & Thomas & Laurent \\\hline
$1,95$m & $1,90$m & $2,01$m & $1,86$m & $1,92$m & $2,03$m & $1,74$m & $1,65$m & $1,97$m \\\hline
$35$pts & $24$pts & $31$pts & $32$pts & $33$pts & $27$pts & $3$pts & $0$pts & $22$pts \\\hline
\end{tabularx}
\end{center}

\begin{enumerate}[font=\color{exer1}]
\item Calcule la taille moyenne des joueurs de cette équipe. Arrondi au cm.
\item Calcul le nombre moyen de points marqués par joueurs, par cette équipe au cours de cette saison.
\end{enumerate}
\end{exer1}

\begin{exer1}
Lors d'une compétition de snowboard, Tom passe deux épreuves : un slalom et une session freestyle en half-pipe.

\begin{enumerate}[font=\color{exer1}]
\item Voici les temps que Tom a réalisés lors de trois descentes de slalom : $2$mn $45$s, $3$mn $1$s et $2$mn $41$s.

\noindent Quel est le temps moyen de Tom sur le slalom ? 

\noindent Pour ce temps Tom obtient $175$ points.
\item Voici maintenant les résultats de Tom sur les trois runs de half-pipe : $187$pts, $236$ pts et $192$pts.

\noindent Quelle est la moyenne des points obtenus par Tom sur cette seconde épreuve?

\item Le score final est la moyenne des points pour le slalom et pour le freestyle. Quel est le score final de Tom ?
\end{enumerate}
\end{exer1}

\begin{exer1}
Voici le discours d'un entraineur de football en fin de saison à son équipe :

\begin{quotation}
\textit{Après avoir marqué $8$ buts lors des $4$ premières rencontres, on au eu un petit passage à vide avec seulement $3$ buts marqués lors des $5$ matchs suivants ! \\Par contre, un grand bravo pour la fin de saison et les $11$ buts marqués sur les $3$ derniers matchs !}
\end{quotation}

Calcule la moyenne de buts marqués par match par l'équipe lors de cette saison.
\end{exer1}

\begin{exer1}
Voici le température en degré Celsius relevées chaque jour d'un mois de novembre : \quad 5 \quad 4 \quad 6 \quad 2 \quad 1 \quad 4 \quad 5 \quad 6 \quad 3 \quad 0 \quad -2 \quad -1 \quad -1 \quad 4 \quad 5 \quad 5 \quad 5 \quad 0 \quad 0 \quad 4 \quad 3 \quad 3 \quad 5 \quad 5 \quad -1 \quad 5 \quad 6 \quad 0 \quad -2 \quad 0

\begin{enumerate}[font=\color{exer1}]
\item Classe les données dans le tableau.
\item Calcule la température moyenne de ce \\mois de novembre (arrondis au dixième).
\end{enumerate}

\vspace{-1.8cm}
\begin{flushright}
\renewcommand{\arraystretch}{1.5}
\begin{tabularx}{9.5cm}{|p{2cm}|*{9}{>{\centering \arraybackslash}X|}}
\hline
\cellcolor{exer1!20}{Température} & $-2$ & $-1$ & $0$ & $1$ & $2$ & $3$ &  $4$ & $5$ & $6$ \\\hline
\cellcolor{exer1!20}{Nb de jours} & & & & & & & & & \\\hline
\end{tabularx}
\end{flushright}
\end{exer1}

\begin{exer2}
Voila, pour le moment, les notes (sur 20) de Bastien au troisième trimestre en mathématiques :

\begin{flushright}
10 \quad 9 \quad 15 \quad 5 \quad 3 \quad 8 \quad 15 \quad 15 \hspace{2cm} \textcolor{white}{.}
\end{flushright}

\vspace{-8mm}
\begin{enumerate}[font=\color{exer2}]
\item Quelle est sa moyenne trimestrielle ?
\item On ajoute une note et sa moyenne augmente. Que peux tu dire sur cette note ?
\item Il obtient $9,5$ à un contrôle. Que se passe-t-il pour sa moyenne ?
\item Le professeur s'est trompé sur deux notes rentrées sur PRONOTE. En fait la moyenne de Bastien est de $12,5$. Effectue les changements.
\end{enumerate}
\end{exer2}

\subsection*{Médiane une série statistique}
\addcontentsline{toc}{subsection}{Médiane une série statistique}

\begin{exer1}
Détermine la médiane de chacune des séries statistiques suivantes.

\begin{enumerate}[font=\color{exer1}]
\item $10$ - $25$ – $32$ – $143$ – $280$ – $301$
\item $8$ – $9.5$ – $14$ – $16$ – $18$ – $19$ – $20$
\item $2$ – $4$ – $7$ – $9$ – $9$ – $10$ – $10$ – $10$ – $11$ – $12$
\item $1$ – $1$ – $1$ – $3$ – $4$ – $4$ – $5$ – $5$ – $5$ – $9$ – $10$
\end{enumerate}
\end{exer1}

\begin{exer1}
Un gérant a relevé le nombre de personnes fréquentant son club sur une semaine.
\begin{center}
\renewcommand{\arraystretch}{1.2}
\begin{tabularx}{\linewidth}{|*{7}{>{\centering \arraybackslash}X|}}
\hline
\rowcolor{exer1!20} Lu & Ma & Me & Je & Ve & Sa & Di \\\hline
$32$ & $38$ & $21$ & $49$ & $60$ & $84$ & $24$\\\hline
\end{tabularx}
\end{center}

\begin{enumerate}[font=\color{exer1}]
\item Calcule le nombre moyen de personnes fréquentant le club par jour.
\item Détermine une médiane de cette série.
\end{enumerate}
\end{exer1}

\begin{exer3}
Sam a relevé les durées des morceaux de sa compilation de rap préférée en min:sec.

\medskip
\noindent $4$:$08$ \quad ; \quad $3$:$19$ \quad ; \quad $4$:$47$ \quad ; \quad $3$:$46$ \quad ; \quad $3$:$15$ \quad ; \quad $3$:$19$ \quad ; \quad $3$:$58$ \quad ; \quad $3$:$50$ \quad ; \quad $3:24$ \quad ; \quad $3$:$55$ \quad ; \quad $3$:$16$ \quad ; \quad $3$:$24$ \quad ; \quad $3$:$07$ 

\begin{enumerate}[font=\color{exer3}]
\item Calcule la durée moyenne des morceaux.
\item Détermine la durée médiane.
\end{enumerate}
\end{exer3}

\subsection*{Représenter une série statistiques}
\addcontentsline{toc}{subsection}{Représenter une série statistiques}

\begin{exer1}
Voici les réponses des élèves d'une classe de 

\vspace{-6mm}
\noindent \begin{minipage}{9.5cm}

\vspace{5mm}
$5^e$ à un sondage portant sur le couleur préféré.

\begin{center}
\renewcommand{\arraystretch}{1.5}
\begin{tabularx}{\linewidth}{|p{2cm}|*{5}{>{\centering \arraybackslash}X|}}
\hline
\cellcolor{exer1!20}{Couleur} & Rouge & Vert & Bleu & Violet & Jaune\\\hline
\cellcolor{exer1!20}{Effectif} & $6$ & $8$ & $5$ & $3$ & $5$ \\\hline
\end{tabularx}
\end{center}

Représente cette série à l'aide d'un diagramme en bâton sur le repère tracé ci-contre.
\end{minipage}
\begin{minipage}{8cm}
\begin{center}
\begin{tikzpicture}[line cap=round,line join=round,>=triangle 45,x=1.0cm,y=1.0cm]
\begin{axis}[x=1cm,y=0.4cm,axis lines=middle,ymajorgrids=true,xmajorgrids=true,
xmin=0,xmax=5.5,ymin=0,ymax=8.5,
xtick={1,2,3,4,5},
xticklabels={Rouge, Vert, Bleu, Violet, Jaune},
ytick={-0,1,...,8},
ylabel=Effectif,]
\clip(0,0) rectangle (5.5,8.6);
%\draw[line width=0.8pt,color=exer1,fill=exer1,fill opacity=0.3] (0.75,0.) rectangle (1.25,6);
%\draw[line width=0.8pt,color=exer1,fill=exer1,fill opacity=0.3] (1.75,0.) rectangle (2.25,8);
%\draw[line width=0.8pt,color=exer1,fill=exer1,fill opacity=0.3] (2.75,0.) rectangle (3.25,5);
%\draw[line width=0.8pt,color=exer1,fill=exer1,fill opacity=0.3] (3.75,0.) rectangle (4.25,3);
%\draw[line width=0.8pt,color=exer1,fill=exer1,fill opacity=0.3] (4.75,0.) rectangle (5.25,5);
\end{axis}
\end{tikzpicture}
\end{center}
\end{minipage}
\end{exer1}

\begin{exer2}
Effectuons maintenant un sondage à l'intérieur de cette classe portant sur votre couleur préféré.

\begin{center}
\renewcommand{\arraystretch}{1.5}
\begin{tabularx}{\linewidth}{|p{2cm}|*{6}{>{\centering \arraybackslash}X|}}
\hline
\cellcolor{exer2!20}{Couleur} & Rouge & Vert & Bleu & Violet & Jaune & Autre\\\hline
\cellcolor{exer2!20}{Effectif} &  &  &  &  &  &  \\\hline
\end{tabularx}
\end{center}

Représente cette série à l'aide d'un diagramme en bâton.
\end{exer2}

\begin{exer2}
La répartition du régime alimentaire d'un sanglier en France est donnée dans le tableau ci-contre. L'objectif est de représenter les données sous \\forme d'un diagramme circulaire.

\begin{enumerate}[font=\color{exer2}]
\item Complète le tableau de proportionnalité \\suivant en écrivant tes calculs.
\item Représente cette répartiton alimentaire à \\l'aide d'une diagramme circulaire
\item Quelle est l'alimentation principale du \\sanglier ?
\item Quels pourcentages représentent les \\graminées et les céréales dans l'alimentation \\du sanglier ?
\end{enumerate}

\vspace{-5.2cm}
\begin{flushright}
\renewcommand{\arraystretch}{1.5}
\begin{tabularx}{9cm}{|p{2.5cm}|*{6}{>{\centering \arraybackslash}X|}}
\hline
\rowcolor{exer2!20} Aliments & Pourcentage & Angles (en degré)\\\hline
\cellcolor{exer2!20}{Fruits forestiers} & $60$ &  \\\hline
\cellcolor{exer2!20}{Graminées} & $18$ &  \\\hline
\cellcolor{exer2!20}{Céréales} & $11$ &  \\\hline
\cellcolor{exer2!20}{Animaux} & $4$ &  \\\hline
\cellcolor{exer2!20}{Racines} & $7$ &  \\\hline
\cellcolor{exer2!20}{Total} & $100$ & $360$ \\\hline
\end{tabularx}
\end{flushright}
\end{exer2}

\begin{exer3}
Madame BIOT et Monsieur BOUTROIS sont tous les deux professeurs de mathématiques et ont tous les deux une classe de cinquième ayant$20$ élèves. Ils souhaitent comparer les notes obtenues par leurs élèves au dernier devoir commun.

\textbf{Note attribuées par Madame Biot : }
7-8-12-12-18-5-11-6-3-8-5-18-9-20-6-16-6-18-7-15

\textbf{Note attribuées par Monsieur Boutrois : }
8-8-9-12-11-8-16-15-7-6-10-10-12-8-10-14-12-11-14-9

\begin{enumerate}[font=\color{exer3}]
\item Construis sur un même graphique les diagrammes en bâtons représentant les deux séries de notes. Utilise deux couleurs
\item Calcule la moyenne de chaque classe
\item En observant le diagramme en bâtons de la question \textcolor{exer3}{\textbf{1.}} et la réponse de la question \textcolor{exer3}{\textbf{2.}}, que pourrait-on conclure ?
\end{enumerate}
\end{exer3}

\end{document}